%! AP Statistics — Unit 5, Lesson 5.2 — Group Activity (Answer Key)
\documentclass[10pt]{article}
\usepackage{apstats-article}
\usepackage{apstats-key}   % pulls in -boxes; do NOT also load -boxes
\newcommand{\TallMath}[1]{$\displaystyle #1\rule[-1.4em]{0pt}{3.2em}$}

\begin{document}

\pageheader{Unit 5, Lesson 5.2}{Group Activity --- Answer Key}
\namepartnerperiod

\vspace{0.15in}

\begin{tcolorbox}[colback=redbg, colframe=black!40, boxrule=0.6pt, arc=2mm,
    left=3mm, right=3mm, top=2mm, bottom=2mm, title=\textbf{Tier R --- Remediate}]
\textbf{Context: Baby weights.} Newborn birth weights in the U.S. are approximately
normal with mean $\mu = 7.5$ lb and standard deviation $\sigma = 1.1$ lb.

\begin{enumerate}
  \item Use the 68--95--99.7 rule to find the range of birth weights for the middle
        95\% of newborns. Show your calculation.
        \ansline{$7.5 \pm 2(1.1) = 7.5 \pm 2.2$}\\
        \ansline{Middle 95\% of newborns: 5.3 lb to 9.7 lb.}
  \item A newborn weighs 6.0 lb. Compute the z-score and determine whether this weight
        is unusually low. Show work.
        \ansline{$z = (6.0 - 7.5)/1.1 \approx -1.36$}\\
        \ansline{About 1.4 SD below the mean --- below average,}\\
        \ansline{but not unusually low (it is within 2 SD of the mean).}
  \item Using the z-score from part (b), the probability of a birth weight \emph{below}
        6.0 lb is approximately 0.086. Interpret this probability in context.
        \ansline{}\\
        \ansline{About 8.6\% of U.S. newborns weigh less than 6.0 lb.}
\end{enumerate}
\end{tcolorbox}

\vspace{0.1in}

\begin{tcolorbox}[colback=redbg, colframe=black!40, boxrule=0.6pt, arc=2mm,
    left=3mm, right=3mm, top=2mm, bottom=2mm, title=\textbf{Tier A --- Approaching Proficiency}]
\textbf{Context: Standardized test.} Scores on a state exam are approximately normal
with $\mu = 72$ and $\sigma = 9$.

\begin{enumerate}
  \item Find $P(X > 85)$. Sketch, compute z-score, and use table/technology.
        \par\vspace{-\parskip}\nointerlineskip\noindent\begin{minipage}[t][1.1in][t]{\linewidth}\vspace*{2pt}\ans{Shade the right tail above 85.}\\
        \ans{$z = (85-72)/9 \approx 1.44$}\\
        \ans{$P(X>85) = P(Z>1.44) = 1 - 0.9251 \approx 0.075$}\end{minipage}
  \item Find $P(60 < X < 80)$. Show all steps.
        \par\vspace{-\parskip}\nointerlineskip\noindent\begin{minipage}[t][1.1in][t]{\linewidth}\vspace*{2pt}\ans{$z_1 = (60-72)/9 \approx -1.33$, \quad $z_2 = (80-72)/9 \approx 0.89$}\\
        \ans{$P(-1.33<Z<0.89) = 0.8133 - 0.0918$}\\
        \ans{$P(60<X<80) \approx 0.72$}\end{minipage}
  \item The top 10\% of scores earn a scholarship. What is the minimum score needed?
        Use inverse normal.
        \par\vspace{-\parskip}\nointerlineskip\noindent\begin{minipage}[t][0.9in][t]{\linewidth}\vspace*{2pt}\ans{Area to the left $= 0.90$, so $z \approx 1.28$}\\
        \ans{$x = 72 + 1.28(9) \approx 83.5$ --- a score of about 84}\end{minipage}
  \item Write a complete sentence interpreting your answer to part (c) in context.
        \ansline{}\\
        \ansline{A student must score at least about 83.5 (84) to earn the scholarship.}
\end{enumerate}
\end{tcolorbox}

\vspace{0.1in}

\begin{tcolorbox}[colback=redbg, colframe=black!40, boxrule=0.6pt, arc=2mm,
    left=3mm, right=3mm, top=2mm, bottom=2mm, title=\textbf{Tier E --- Extension}]
\textbf{Context: Daily rainfall.} A weather model predicts that daily rainfall in a
city follows $N(1.2\text{ cm}, 0.4\text{ cm})$ during summer.

\begin{enumerate}
  \item Find the probability that a randomly selected summer day has rainfall between
        0.8 cm and 2.0 cm.
        \par\vspace{-\parskip}\nointerlineskip\noindent\begin{minipage}[t][0.9in][t]{\linewidth}\vspace*{2pt}\ans{$z_1 = (0.8-1.2)/0.4 = -1$, \quad $z_2 = (2.0-1.2)/0.4 = 2$}\\
        \ans{$P(-1<Z<2) = 0.9772 - 0.1587 \approx 0.819$}\end{minipage}
  \item Find the 5th and 95th percentiles. What do these values mean in context?
        \par\vspace{-\parskip}\nointerlineskip\noindent\begin{minipage}[t][0.9in][t]{\linewidth}\vspace*{2pt}\ans{$z = \pm 1.645$: \quad $1.2 - 1.645(0.4) \approx 0.54$ cm, \quad $1.2 + 1.645(0.4) \approx 1.86$ cm}\\
        \ans{On 5\% of summer days rainfall is below 0.54 cm;}\\
        \ans{on 5\% of summer days it is above 1.86 cm.}\end{minipage}
  \item The most extreme 2\% of rainfall days (highest and lowest) trigger weather
        alerts. Find the boundary values. (\textit{Hint:} split 2\% into two tails ---
        VAR-6.B.2d.)
        \par\vspace{-\parskip}\nointerlineskip\noindent\begin{minipage}[t][0.9in][t]{\linewidth}\vspace*{2pt}\ans{1\% in each tail: $z = \pm 2.326$}\\
        \ans{Lower: $1.2 - 2.326(0.4) \approx 0.27$ cm; \quad upper: $1.2 + 2.326(0.4) \approx 2.13$ cm}\\
        \ans{Alerts for days below 0.27 cm or above 2.13 cm.}\end{minipage}
  \item Is the normal model appropriate for daily rainfall? What features of the data
        would you check before applying this model?
        \ansline{}\\
        \ansline{Probably not: rainfall is right-skewed and can't go below 0 cm.}\\
        \ansline{Check a histogram or normal probability plot for symmetry/bell shape.}
\end{enumerate}
\end{tcolorbox}

\end{document}
